% \section{Total Migration: Case Studies}
%     - The migration pipeline: target discovery → inference → injection
%       → mix compile in original project context.
%     - 2-3 selected projects: report migration rate
%       (compile-pass + no-type-warning / total untyped functions).
%     - Per-project failure breakdown:
%         a. Predicted type triggers a type warning
%         b. Predicted type fails to compile (syntax / unresolved reference)
%         c. Predicted type compiles cleanly
%     - Compare intrinsic distance (4.5) vs extrinsic compile-pass rate.
%       Does the distance metric predict migration success?

% \section{Discussion}
%     - What worked, what didn't.
%     - Where the model's failures concentrate (structs, dep types, atom-vs-binary).
%     - Limitations: held-out projects vs in-project, dataset size,
%       no dependency-resolution at training time.

% \section{Summary}
%     - Reframe: checker-verified supervision produced two models; the
%       no-gradual model reached 65.2% safe-and-precise typecheck vs the
%       both-pass model's 38.3% (Subsection~\ref{sec:extrinsic-eval}),
%       and X% end-to-end migration on case-study projects; the pipeline is
%       demonstrably functional.

The intrinsic evaluation of Section~\ref{sec:compatibility-eval} measures prediction quality against verified reference annotations.
Total migration poses the complementary question: for functions that carry \emph{no} annotation at all---and therefore no reference---can the model supply annotations that the typechecker accepts?
This chapter defines the migration task, introduces the \emph{cascading weakening strategy} used to guarantee coverage, and reports case studies on real projects.

\section{Task and Setting}
\label{sec:migration-task}

A \emph{migration target} is a function definition in the corpus with no associated \texttt{@spec}.
Such functions are invisible to the translation pipeline of Chapter~\ref{ch:implementation}: with no TypeSpec to translate, they remain unannotated after translation, and Dialyzer accepts them vacuously.
In the projects studied, migration targets constitute the majority of all function definitions.

The migration state is what we term \emph{State A}: each target project retains the translated \texttt{@assert\_type\_form} annotations produced by the prototype for every function that had a \texttt{@spec}, and the model supplies annotations only for the remaining untyped functions.
This mirrors the intended deployment: existing developer-written specifications are translated mechanically; the model fills the gaps.
Total migration is achieved when every function in the project carries an annotation accepted by the new typechecker.

Because no reference type exists for a migration target, the reference-based metric of Section~\ref{sec:compatibility-eval} does not apply.
The only available correctness signal is the typechecker itself, invoked through \texttt{mix compile} in the original project context, where dependencies resolve and remote types are visible.
We therefore adopt the typechecker as the migration oracle:
\[
\mathsf{accept}(f, \tau) \iff \text{\texttt{mix compile} emits no type warning and no compile error for } f \text{ annotated with } \tau.
\]

\section{The Cascading Weakening Strategy}
\label{sec:cascade}

A single-shot prediction pipeline would leave every rejected prediction unmigrated.
Instead, we exploit the structure of the set-theoretic type system: any annotation can be \emph{weakened}---moved upward in the precision order---by deterministic syntactic transformations, and the maximally weakened annotation is accepted unconditionally.
The strategy is therefore to generate once at full precision and, on rejection, apply successive weakenings until the oracle accepts.
Crucially, the model is invoked exactly once per target; all subsequent levels are rewrites of its output, not regenerations.

\paragraph{Cascade levels.}
Let $\tau_0$ be the raw model output for a target $f$.
The cascade defines a sequence of candidate annotations $\tau_0, \tau_1, \tau_2, \tau_3$:

\begin{description}
    \item[Level 0 (full precision).]
    $\tau_0$ as generated, after syntactic validation.
    If $\tau_0$ is malformed (unbalanced brackets, truncated generation), a repair pass truncates the type to its last syntactically complete top-level union arm; if no complete arm exists, the entry skips directly to Level~3.

    \item[Level 1 (struct compaction).]
    $\tau_1 = \mathcal{C}(\tau_0)$, where the compaction operator $\mathcal{C}$ rewrites every struct-shaped map to its compact form, retaining only the struct tag:
    \[
    \mathcal{C}\bigl(\texttt{\%\{}k_1 \Rightarrow t_1, \ldots, \texttt{:\_\_struct\_\_} \Rightarrow \texttt{:"M"}\texttt{\}}\bigr)
    \;=\;
    \texttt{\%\{:\_\_struct\_\_} \Rightarrow \texttt{:"M"\}}.
    \]
    Compaction removes the field-level commitments that are the dominant source of struct-shape errors (Section~\ref{sec:qualitative}), including hallucinated field sets, while preserving the nominal identity of the struct.

    \item[Level 2 (positional gradualization).]
    $\tau_2 = \mathcal{G}(\tau_1)$, where the gradualization operator $\mathcal{G}$ replaces each input position of the function type with \texttt{dynamic()} and widens the output to the union of the output with \texttt{dynamic()}:
    \[
    \mathcal{G}\bigl(\texttt{(} \iota_1, \ldots, \iota_n \rightarrow \omega \texttt{)}\bigr)
    \;=\;
    \texttt{(dynamic()}, \ldots, \texttt{dynamic()} \rightarrow \omega\texttt{)}.
    \]
    Inputs are replaced with \texttt{dynamic()} rather than \texttt{term()} deliberately: widening an input to \texttt{term()} strengthens the obligation on the function body (it must now be verified against arbitrary inputs), whereas \texttt{dynamic()} discharges the obligation by consistency, which is the intended role of the gradual type.
    The output $\omega$ is retained at this level, preserving the model's return-shape prediction, which the intrinsic evaluation shows to be the more reliable component.

    \item[Level 3 (full gradual fallback).]
    $\tau_3 = \texttt{(dynamic()}, \ldots, \texttt{dynamic()} \rightarrow \texttt{dynamic())}$.
    This annotation is accepted unconditionally, guaranteeing that the cascade terminates with every target migrated.
    Entries resolved at Level~3 carry no information from the model; they are counted as coverage without precision.
\end{description}

Each level is a deterministic function of the previous one, so the cascade requires no additional model inference and no randomness; the entire migration of a project is reproducible from the Level-0 predictions alone.

\paragraph{Batched execution.}
Because the oracle operates at project granularity (one \texttt{mix compile} verifies all annotations simultaneously), the cascade is executed in rounds rather than per function:
all targets are annotated at Level~0 and the project is compiled once; the warning and error output is parsed to attribute failures to individual functions; failed targets are rewritten to Level~1 and the project recompiled; and so on.
A full migration therefore costs at most four compilations per project, independent of the number of targets.
Algorithm~\ref{alg:cascade} summarises the procedure.

\begin{algorithm}
\caption{Cascading migration of a project $P$}
\label{alg:cascade}
\begin{algorithmic}[1]
\State $T \gets$ functions of $P$ without \texttt{@spec} \Comment{target discovery}
\State $\mathit{pred} \gets$ model predictions for all $f \in T$ \Comment{single inference pass}
\State $\mathit{level}[f] \gets 0$ for all $f \in T$
\For{$\ell = 0$ \textbf{to} $3$}
    \State annotate each unresolved $f \in T$ with $\tau_\ell(f)$; \texttt{mix compile} $P$
    \State $W \gets$ functions attributed a type warning or compile error
    \For{$f \in T \setminus W$ unresolved}
        \State mark $f$ resolved at level $\ell$
    \EndFor
\EndFor
\State \Return level distribution $\{\,f \mapsto \mathit{level}[f]\,\}$
\end{algorithmic}
\end{algorithm}

\section{Outcome Classification and Metrics}
\label{sec:migration-metrics}

The primary result of a migration run is not a single pass rate but the \emph{level distribution}: the fraction of targets resolved at each cascade level.
The distribution refines a binary migration rate into a precision profile:
\begin{itemize}
    \item \textbf{L0 rate}: targets migrated at full model precision---the strongest outcome;
    \item \textbf{L1 rate}: targets whose field-level struct commitments failed but whose nominal and shape structure was accepted;
    \item \textbf{L2 rate}: targets for which only the predicted return shape survived;
    \item \textbf{L3 rate}: targets migrated with no usable model contribution.
\end{itemize}
The cumulative migration rate at every level is 100\% by construction; the informative quantity is how much of that coverage is achieved at high precision.
We additionally record, at each level, the split between the two failure modes---type warnings (the annotation is well-formed but inconsistent with the body) and compile errors (the annotation is syntactically or referentially invalid)---since the former reflects model imprecision and the latter reflects generation pathology.

Two annotations are reported per track: the same cascade is run once with the Track~1 (precision-trained) adapter and once with the Track~2 (coverage-trained) adapter, allowing the precision--coverage trade-off observed intrinsically in Section~\ref{sec:two-track} to be re-examined under deployment conditions.
The hypothesis suggested by the intrinsic results is that Track~1 resolves more targets at Level~0 on precisely-typed code but pushes more targets to Levels~2--3, while Track~2 exhibits a flatter distribution.

\section{Case Studies}
\label{sec:migration-case-studies}

The cascade was executed on \emph{N} projects selected to span the difficulty spectrum identified in the per-subproject analysis:
one standard-library-style codebase (\texttt{elixir-lang/X}), one framework-heavy codebase (\texttt{ScenicFramework/Scenic}), and one application codebase (\texttt{Y}).
Table~\ref{tab:migration-results} reports the level distributions.

\begin{table}[h]
\centering
\begin{tabular}{llrrrrrr}
\toprule
\textbf{Project} & \textbf{Track} & \textbf{Targets} & \textbf{L0} & \textbf{L1} & \textbf{L2} & \textbf{L3} & \textbf{L0+L1} \\
\midrule
\texttt{project-A} & 1 & X & X\% & X\% & X\% & X\% & X\% \\
                   & 2 & X & X\% & X\% & X\% & X\% & X\% \\
\texttt{project-B} & 1 & X & X\% & X\% & X\% & X\% & X\% \\
                   & 2 & X & X\% & X\% & X\% & X\% & X\% \\
\texttt{project-C} & 1 & X & X\% & X\% & X\% & X\% & X\% \\
                   & 2 & X & X\% & X\% & X\% & X\% & X\% \\
\bottomrule
\end{tabular}
\caption{Cascade level distribution per project and track. L0+L1 aggregates the outcomes in which the model's structural prediction survives.}
\label{tab:migration-results}
\end{table}

% Placeholder subsections to fill after the runs:
% \paragraph{Failure-mode breakdown.} warnings vs compile errors per level.
% \paragraph{Track comparison.} does the intrinsic precision/coverage trade-off
%   reproduce under the compiler oracle?
% \paragraph{Examples.} one L0, one L1-recovered, one L3 case with commentary.

\section{Relation to the Intrinsic Metric}
\label{sec:migration-vs-intrinsic}

The cascade oracle and the reference-based metric of Section~\ref{sec:compatibility-eval} measure different properties, and their disagreement is informative.
An annotation strictly wider than the function's true type is imprecise under the reference metric yet accepted by the checker; conversely, an annotation more precise than a loose reference may be rejected.
For the subset of test-set functions with references, we re-ran the cascade (with the typechecker, not the reference, as oracle) and cross-tabulated cascade level against the tier assigned by the subtype metric.
% Table placeholder: tier (rows) x resolved level (columns).
A strong concentration of tier-1/tier-2 predictions at Level~0 would validate the intrinsic metric as a proxy for deployability; systematic divergence would delimit what reference-based evaluation can and cannot predict about migration outcomes.
